\documentclass[11pt,a4paper]{article}
\usepackage[T2A]{fontenc}
\usepackage[utf8]{inputenc}
\usepackage[bulgarian]{babel}
\usepackage[margin=2cm]{geometry}
\usepackage{cmap}
\usepackage{booktabs,tabularx,xltabular,array,multirow}
\usepackage[labelfont=bf,labelsep=colon,font=small]{caption}
\usepackage{parskip}
\usepackage{float}
\usepackage{enumitem}
\usepackage{eurosym}
\usepackage{xcolor}
\usepackage{tikz}
\usetikzlibrary{positioning,arrows.meta,fit,calc}
\usepackage{pgfgantt}
\usepackage{pgfplots}
\pgfplotsset{compat=1.18}
\usepackage[hidelinks]{hyperref}
\definecolor{wpfill}{HTML}{F2D46B}
\definecolor{wpdraw}{HTML}{B8901A}
\definecolor{tkfill}{HTML}{FFF3B0}
\definecolor{tkdraw}{HTML}{C9A227}
\definecolor{msfill}{HTML}{8C2F1B}
\definecolor{mtfill}{HTML}{FFF9DD}
\setlist{nosep,leftmargin=1.2em}
\renewcommand{\arraystretch}{1.12}
\newcolumntype{R}[1]{>{\raggedleft\arraybackslash}p{#1}}
\newcolumntype{C}[1]{>{\centering\arraybackslash}p{#1}}
\newcolumntype{L}[1]{>{\raggedright\arraybackslash}p{#1}}
\setlength{\LTpre}{6pt}\setlength{\LTpost}{6pt}
\begin{document}
\begin{center}
{\Large\bfseries YTKI: Интелигентна мултиагентна система\\ за мониторинг и реакция при инциденти в облачни услуги\par}
\vspace{4pt}
{\small Домашна работа 3: ресурси, график и бюджет на работната програма\par}
\vspace{8pt}
\begin{tabular}{ll}
\toprule
\textbf{Име и фамилия} & \textbf{Факултетен номер}\\
\midrule
Йоан Байчев & 0MI0600328\\
Тугай Кафа & 4MI0600275\\
Калоян Марков & 4MI0600392\\
Илиан Запрянов & 8MI0600405\\
\bottomrule
\end{tabular}
\end{center}

\textbf{Тема:} Интелигентна система за мониторинг, базирана на AI агенти. Системата работи 24/7 и в реално време приема алерти, логове и health сигнали от различни landscape-и и среди (dev, staging, int-prod, prod). Тя класифицира инцидентите и помага на DevOps, SRE и operations екипите да ги разследват и решават по-бързо.


\section*{2\quad Описание на работната програма по проекта}
\addcontentsline{toc}{section}{2 Описание на работната програма}
В Домашна работа 2 определихме целта на проекта, шестте подцели с междинните резултати MS1--MS6 и разделихме работата на 8 работни пакета и 38 задачи (WBS, фиг.~1 там). Тук описваме как точно ще се свърши тази работа: логиката на работната програма, съдържанието на всяка задача, разпределението ѝ във времето, резултатите и ресурсите, които са нужни за нея. Разделите и таблиците продължават номерацията от Домашна работа 2: подробните таблици 2.1--2.8 заменят съкратените, в които задачите бяха само изброени.

\subsection*{2.2\quad Работна програма и задачи}

\textbf{Обща методология на работната програма.} Работната програма следва пътя на един алерт през системата (фиг.~2). Всеки технически пакет изгражда един или два слоя от архитектурата в Домашна работа 1 и завършва с междинен резултат, който може да се провери. Така всяка подцел има собственик и ясен край, а напредъкът се вижда през работещ софтуер, а не през отчети.
\begin{itemize}
\item \textbf{Първо знание, после код.} РП1 (месеци 1--2) измерва базовите показатели, събира изискванията и определя таксономията и архитектурата. Без него нямаме спрямо какво да оценим ефекта и какво да класифицираме.
\item \textbf{Данните преди агентите.} РП2 изгражда единния поток от алерти, конекторите и графа на зависимостите. Агентите в РП3 и моделите в РП4 работят само с данни, които вече са нормализирани и маскирани.
\item \textbf{Разследване и класификация паралелно.} РП3 и РП4 вървят едновременно (месеци 3--7): докато се правят агентите, се анотират историческите инциденти, а докладът на агентите (3.3) е вход за класификацията (4.2).
\item \textbf{Реакцията накрая.} РП5 автоматизира действията едва когато има анализ и категория, на които може да се разчита. Опасните действия (рестарт, автоматично затваряне) винаги минават през одобрени правила или през човек.
\item \textbf{Prod последна.} РП6 пуска системата в prod само след като тестовете (6.1) и проверката за сигурност (6.2) са минали. Пилотът (6.5) измерва ефекта спрямо базовите стойности от 1.1, а обратната връзка от него се връща в преобучението (4.4).
\item \textbf{Хоризонтални пакети.} РП7 (комуникации) и РП8 (управление) текат през целия проект и свързват техническата работа със заинтересованите лица.
\end{itemize}


\setcounter{figure}{1}
\begin{figure}[ht]
\centering
\begin{tikzpicture}[
  wp/.style={draw=wpdraw, fill=tkfill, rounded corners=2pt, align=center, font=\scriptsize, minimum height=1.75cm, text width=2.62cm, inner sep=3pt},
  ar/.style={-{Stealth[length=2mm]}, draw=black!70, thick},
  lb/.style={font=\tiny, align=center, text=black!75}]
\node[wp] (w1) at (0,0) {\textbf{РП1} Анализ и\\проектиране\\м. 1--2, MS1\\[2pt]\tiny\itshape изисквания, архитектура, базови стойности};
\node[wp] (w2) at (3.45,0) {\textbf{РП2} Приемане и\\разпространение\\м. 2--4, MS2\\[2pt]\tiny\itshape нормализирани и маскирани данни};
\node[wp] (w3) at (6.9,1.3) {\textbf{РП3}\\AI агенти за\\разследване\\м. 3--6, MS3\\[2pt]\tiny\itshape доклад и хипотеза за причината};
\node[wp] (w4) at (6.9,-1.3) {\textbf{РП4}\\Класификация и\\самообучение\\м. 3--7, MS4\\[2pt]\tiny\itshape категория с вероятност};
\node[wp] (w5) at (10.35,0) {\textbf{РП5} Оперативна\\реакция\\м. 5--7, MS5\\[2pt]\tiny\itshape насочване, тикети,\\проверки};
\node[wp] (w6) at (13.8,0) {\textbf{РП6} Тестване,\\внедряване, оценка\\м. 6--9, MS6\\[2pt]\tiny\itshape работеща система и измерен ефект};
\draw[ar] (w1) -- (w2);
\draw[ar] (w2.east) -- ++(0.35,0) |- (w3.west);
\draw[ar] (w2.east) -- ++(0.35,0) |- (w4.west);
\draw[ar] (w3) -- (w4);
\draw[ar] (w3.east) -- ++(0.35,0) |- (w5.west);
\draw[ar] (w4.east) -- ++(0.35,0) |- (w5.west);
\draw[ar] (w5) -- (w6);
\draw[ar, dashed] (w6.south) |- (10.35,-2.75) -- (7.5,-2.75) node[lb, pos=0.5, below]{обратна връзка от пилота и преобучение (4.4)} -- ([xshift=6mm]w4.south);
\node[draw=wpdraw, fill=wpfill, rounded corners=2pt, font=\scriptsize, minimum height=0.45cm, minimum width=16.3cm, anchor=north] at (6.9,-3.2) {\textbf{РП7} Комуникации \qquad \textbf{РП8} Управление на проекта \qquad (м. 1--9, през целия проект)};
\end{tikzpicture}
\caption{Логика на работната програма: пакетите следват пътя на алерта през слоевете на YTKI; под всеки пакет е резултатът, който той дава на следващите. РП7 и РП8 обхващат целия проект.}
\label{fig:logic}
\end{figure}


\textbf{Как структурата води до целта.} Всеки от трите индикатора от Домашна работа 2 зависи от конкретни пакети. По-ниският MTTR (поне $-30\%$) идва от готовия анализ на агентите (РП3) и от това, че тикетът стига направо до правилния екип (РП5). По-малкото алерти за ръчна обработка (поне $-40\%$) идват от отсяването на фалшивите и дублираните алерти (2.1, РП4) и от автоматичното им затваряне (5.1). Наличността от 99,5\% се осигурява в РП6 чрез тестове за натоварване, поетапно внедряване и наблюдение на самата YTKI. Ефектът се измерва в 6.5 спрямо базовите стойности от 1.1, така че първата и последната задача на проекта са свързани.

\textbf{Организация на работата.} Проектът е разделен на 18 двуседмични спринта (два на месец), които са и единицата в графика (фиг.~3). Всяка задача има водеща роля, която отговаря за резултата ѝ (табл.~3). В проекта участват 14 души от организацията, но повечето не са заети изцяло: средно натоварването е 4,75 души на пълен работен ден, а най-голямо е в месец 6 (около 6,5). Затова ресурсите са планирани в човекомесеци (ч.м.) за всяка задача, а не по брой хора.

\textbf{Критичен път и резерви.} Най-дългата верига от зависими задачи е 1.5 $\rightarrow$ 2.1 $\rightarrow$ 2.2 $\rightarrow$ 3.1 $\rightarrow$ 3.2 $\rightarrow$ 3.3 $\rightarrow$ 4.2 $\rightarrow$ 4.5 $\rightarrow$ 6.1/6.2 $\rightarrow$ 6.3 $\rightarrow$ 6.5. Забавяне по нея измества пилота, затова тези задачи се следят с приоритет в 8.3. Задачите извън нея (напр. 2.4, 3.4, 5.4, 7.x) имат свобода от един до няколко спринта. Резервът за непредвидени разходи е 10\% от бюджета (раздел~3).

\textbf{Описание на работните пакети.} Таблици 2.1--2.8 описват всеки пакет по шаблона. За всяка задача са дадени: цел и същност (какво и как се прави), очакван резултат (deliverable) с месеца, в който се представя, връзки с други задачи и необходимите ресурси: екип в човекомесеци по роли, материални ресурси (инфраструктура, услуги, лицензи) и нематериални ресурси (достъпи, данни, знание, време на хора извън екипа). Последният ред на всяка задача дава разходите ѝ в евро; как са изчислени, е обяснено в раздел~3.

\renewcommand{\thetable}{2.\arabic{table}}\setcounter{table}{0}
{\small
\begin{xltabular}{\textwidth}{@{}L{2.55cm}X@{}}
\caption{Описание на работен пакет РП1}\\
\toprule
\multicolumn{2}{@{}p{\textwidth}@{}}{\textbf{Номер и име на работния пакет:} РП1: Анализ и проектиране}\\*
\multicolumn{2}{@{}l@{}}{\textbf{Начало (месец):} 1\qquad \textbf{Край (месец):} 2}\\*
\multicolumn{2}{@{}p{\textwidth}@{}}{\textbf{Цел:} (подцел 1) Анализ на текущия процес по обработка на алерти, измерване на базовите показатели и определяне на изискванията и архитектурата на YTKI}\\*
\multicolumn{2}{@{}p{\textwidth}@{}}{\textbf{Общо ресурси:} 5,75 ч.м.; 37~800~\euro{} (персонал 34~700~\euro{}, други преки разходи 3~100~\euro{})}\\
\midrule
\endfirsthead
\multicolumn{2}{@{}l@{}}{\textit{Таблица 2.\arabic{table}: РП1 (продължение)}}\\ \midrule
\endhead
\bottomrule
\endlastfoot
\multicolumn{2}{@{}p{\textwidth}@{}}{\textbf{Задача 1.1: Анализ на текущия процес и базовите показатели}\hfill месец 1}\\*
\textit{Цел и същност} & Наблюдаваме няколко дежурни смени и интервюираме дежурните инженери, за да опишем как сега минава един алерт от получаването до затварянето. От Splunk и тикет системата изтегляме историята на алертите и инцидентите за последните 6 месеца и изчисляваме базовите стойности на MTTD, MTTR, броя алерти на смяна и дела на фалшивите и дублираните алерти.\\*
\textit{Резултат} & D1.1 Доклад за текущото състояние и базовите показатели (месец 1)\\*
\textit{Връзки} & Дава входните данни за 1.2 и 1.3. Базовите стойности са еталонът, спрямо който в 6.5 се оценява ефектът на YTKI.\\*
\textit{Екип (ч.м.)} & бизнес анализатор 0,5; SRE инженер 0,5; data scientist 0,5 (общо 1,5 ч.м.)\\*
\textit{Материални ресурси} & Наличната среда за анализ на данни (Jupyter). Няма нови разходи.\\*
\textit{Нематериални ресурси} & Достъп само за четене до Splunk, тикет системата и графика на дежурствата; около 10 интервюта по 1 час с дежурни инженери.\\*
\textit{Разходи} & персонал \textbf{8~400~\euro{}} (няма други преки разходи)\\
\midrule[0.3pt]
\multicolumn{2}{@{}p{\textwidth}@{}}{\textbf{Задача 1.2: Изисквания на заинтересованите лица}\hfill месец 2}\\*
\textit{Цел и същност} & Провеждаме работни срещи с групите от анализа в Домашна работа 1: SRE и дежурни, собственици на услуги, ръководство и екипа по сигурност. Изискванията записваме като потребителски истории с критерии за приемане и ги подреждаме по приоритет (MoSCoW).\\*
\textit{Резултат} & D1.2 Спецификация на изискванията (месец 2)\\*
\textit{Връзки} & Използва D1.1. Вход за архитектурата (1.5) и за backlog-а на РП2–РП5; критериите за приемане стават тестови сценарии в 6.1.\\*
\textit{Екип (ч.м.)} & бизнес анализатор 0,75; софтуерен архитект 0,25; ръководител на проекта 0,25 (общо 1,25 ч.м.)\\*
\textit{Материални ресурси} & Зала и кетъринг за три полудневни работни срещи (300~\euro{}); наличната Jira за изискванията.\\*
\textit{Нематериални ресурси} & Време на представителите на заинтересованите лица (около 15 души по 4 часа).\\*
\textit{Разходи} & персонал 6~850~\euro{} + други 300~\euro{} = \textbf{7~150~\euro{}}\\
\midrule[0.3pt]
\multicolumn{2}{@{}p{\textwidth}@{}}{\textbf{Задача 1.3: Таксономия на инцидентите и критерии за класификация}\hfill месец 1–2}\\*
\textit{Цел и същност} & Заедно с опитни SRE инженери определяме измерими критерии за трите категории (фалшив алерт, деградация, прекъсване): засегнати SLO, скорост на изразходване на error budget, дял на засегнатите заявки. За всяка категория подбираме реални примери от историята.\\*
\textit{Резултат} & D1.3 Таксономия на инцидентите с критерии и примери (месец 2)\\*
\textit{Връзки} & Използва данните от 1.1. Основа за анотирането в 4.1, за промптовете в 4.2 и за правилата за ескалация в 5.1.\\*
\textit{Екип (ч.м.)} & SRE инженер 0,5; data scientist 0,25 (общо 0,75 ч.м.)\\*
\textit{Материални ресурси} & Няма.\\*
\textit{Нематериални ресурси} & Експертното знание на SRE екипа; исторически инциденти и postmortem документи; дефинициите на SLO на услугите.\\*
\textit{Разходи} & персонал \textbf{4~500~\euro{}} (няма други преки разходи)\\
\midrule[0.3pt]
\multicolumn{2}{@{}p{\textwidth}@{}}{\textbf{Задача 1.4: Изисквания за сигурност и съответствие (GDPR, NIS2)}\hfill месец 2}\\*
\textit{Цел и същност} & Логовете съдържат лични данни, затова правим оценка на въздействието върху защитата на данните (DPIA). Определяме кои данни може да стигат до LLM и при какви условия (обработка в ЕС, без съхранение на заявките от доставчика), кой до какво има достъп, както и сроковете и съдържанието на докладите по NIS2.\\*
\textit{Резултат} & D1.4 Политика за сигурност и защита на данните, вкл. правила за данните към LLM (месец 2)\\*
\textit{Връзки} & Определя изискванията към маскирането (2.5), защитните механизми (3.5), докладите по NIS2 (5.5) и проверката за сигурност (6.2).\\*
\textit{Екип (ч.м.)} & специалист по сигурност 0,75 (общо 0,75 ч.м.)\\*
\textit{Материални ресурси} & Външна правна консултация по GDPR и NIS2 (2 000~\euro{}).\\*
\textit{Нематериални ресурси} & Участие на длъжностното лице по защита на данните (DPO) на организацията; договорните условия на доставчика на LLM.\\*
\textit{Разходи} & персонал 4~875~\euro{} + други 2~000~\euro{} = \textbf{6~875~\euro{}}\\
\midrule[0.3pt]
\multicolumn{2}{@{}p{\textwidth}@{}}{\textbf{Задача 1.5: Архитектура на системата и интерфейси на агентите}\hfill месец 1–2}\\*
\textit{Цел и същност} & Детайлизираме слоевете от фиг. 1 в Домашна работа 1: технологии, формат на съобщенията, договорите (JSON схеми) между агентите и мястото на всяка услуга в облака. LLM модела избираме след кратко сравнение на 2–3 модела върху реални алерти. Важните решения записваме като architecture decision records (ADR).\\*
\textit{Резултат} & D1.5 Архитектурен документ на YTKI (месец 2)\\*
\textit{Връзки} & Използва D1.2–D1.4. По него се изграждат РП2–РП5; заедно с D1.2 и D1.4 формира MS1.\\*
\textit{Екип (ч.м.)} & софтуерен архитект 1; AI/LLM инженер 0,25; платформен инженер 0,25 (общо 1,5 ч.м.)\\*
\textit{Материални ресурси} & Изолирана облачна среда за проби (500~\euro{}); LLM API за сравнението на моделите (300~\euro{}).\\*
\textit{Нематериални ресурси} & Съществуващият прототип като отправна точка; опитът на екипа с облачната платформа.\\*
\textit{Разходи} & персонал 10~075~\euro{} + други 800~\euro{} = \textbf{10~875~\euro{}}\\
\end{xltabular}}
{\small
\begin{xltabular}{\textwidth}{@{}L{2.55cm}X@{}}
\caption{Описание на работен пакет РП2}\\
\toprule
\multicolumn{2}{@{}p{\textwidth}@{}}{\textbf{Номер и име на работния пакет:} РП2: Приемане и разпространение на данните}\\*
\multicolumn{2}{@{}l@{}}{\textbf{Начало (месец):} 2\qquad \textbf{Край (месец):} 4}\\*
\multicolumn{2}{@{}p{\textwidth}@{}}{\textbf{Цел:} (подцел 2) Изграждане на единен поток в реално време, през който алертите, логовете и health сигналите от всички среди постъпват нормализирани}\\*
\multicolumn{2}{@{}p{\textwidth}@{}}{\textbf{Общо ресурси:} 5,75 ч.м.; 40~700~\euro{} (персонал 33~500~\euro{}, други преки разходи 7~200~\euro{})}\\
\midrule
\endfirsthead
\multicolumn{2}{@{}l@{}}{\textit{Таблица 2.\arabic{table}: РП2 (продължение)}}\\ \midrule
\endhead
\bottomrule
\endlastfoot
\multicolumn{2}{@{}p{\textwidth}@{}}{\textbf{Задача 2.1: Ingestion слой: прокси, парсване и нормализиране}\hfill месец 2–3}\\*
\textit{Цел и същност} & Изграждаме прокси, което приема алертите по webhook от Prometheus Alertmanager, Cloud Monitoring и Splunk. Всеки алерт се парсва и превежда към единен модел (по семантиката на OpenTelemetry), а дубликатите се отсяват по отпечатък (fingerprint). Тук изграждаме и общата dev/staging среда на проекта.\\*
\textit{Резултат} & D2.1 Ingestion услуга с единен модел на алерта (месец 3)\\*
\textit{Връзки} & Следва архитектурата (1.5). Предпоставка за 2.2 и за агентите в РП3; нормализираните алерти се ползват и в 4.1.\\*
\textit{Екип (ч.м.)} & платформен инженер 1; DevOps/интеграционен инженер 0,5 (общо 1,5 ч.м.)\\*
\textit{Материални ресурси} & Облачна среда dev/staging за целия проект (Kubernetes клъстер, опашка за съобщения, база данни): 8 месеца по около 900~\euro{} (7 200~\euro{}).\\*
\textit{Нематериални ресурси} & Права за настройка на webhook-ове в инструментите за мониторинг на четирите среди.\\*
\textit{Разходи} & персонал 8~700~\euro{} + други 7~200~\euro{} = \textbf{15~900~\euro{}}\\
\midrule[0.3pt]
\multicolumn{2}{@{}p{\textwidth}@{}}{\textbf{Задача 2.2: Broadcasting слой: абонаментен модел и споделен контекст}\hfill месец 2–3}\\*
\textit{Цел и същност} & Нормализираните алерти се публикуват в теми (topics), а агентите и таблото се абонират за тях. За всеки инцидент има общ обект с контекст, който всеки абонат допълва. Схемата е версионирана, за да не се чупят абонатите при промени.\\*
\textit{Резултат} & D2.2 Broadcasting услуга със споделен контекст (месец 3)\\*
\textit{Връзки} & Използва D2.1. От нея четат Planning агентът (3.1), агентите в 3.2 и таблото (5.4).\\*
\textit{Екип (ч.м.)} & платформен инженер 0,75; DevOps/интеграционен инженер 0,25 (общо 1 ч.м.)\\*
\textit{Материални ресурси} & Управлявана pub/sub услуга и кеш в общата облачна среда от 2.1. Няма нови разходи.\\*
\textit{Нематериални ресурси} & Насоките на платформения екип на организацията за работа с опашки.\\*
\textit{Разходи} & персонал \textbf{5~800~\euro{}} (няма други преки разходи)\\
\midrule[0.3pt]
\multicolumn{2}{@{}p{\textwidth}@{}}{\textbf{Задача 2.3: Конектори към логове, метрики и health endpoint-и}\hfill месец 3–4}\\*
\textit{Цел и същност} & Пишем общ слой от конектори към Cloud Logging, Prometheus, OpenSearch и health endpoint-ите на услугите. Всички имат еднакъв интерфейс, ограничение на броя заявки и само права за четене, за да могат агентите безопасно да ги ползват като инструменти (tools).\\*
\textit{Резултат} & D2.3 Библиотека от конектори към източниците на данни (месец 4)\\*
\textit{Връзки} & Агентите в 3.2 ги използват като инструменти; използват ги и проверките в 5.3.\\*
\textit{Екип (ч.м.)} & платформен инженер 1; DevOps/интеграционен инженер 0,5 (общо 1,5 ч.м.)\\*
\textit{Материални ресурси} & Няма нови разходи (работи в общата среда).\\*
\textit{Нематериални ресурси} & Сервизни акаунти с права само за четене във всичките четири среди, вкл. prod (с одобрение от екипа по сигурност); документацията на API-тата.\\*
\textit{Разходи} & персонал \textbf{8~700~\euro{}} (няма други преки разходи)\\
\midrule[0.3pt]
\multicolumn{2}{@{}p{\textwidth}@{}}{\textbf{Задача 2.4: Граф на зависимостите и мапинг на компонентите}\hfill месец 2–4}\\*
\textit{Цел и същност} & Разширяваме графа от прототипа с данни от трасировките (OpenTelemetry) и от наличния каталог на услугите. За всяка услуга записваме отговорния екип и дежурната ротация и ги потвърждаваме със собствениците, докато покритието стане поне 95\%.\\*
\textit{Резултат} & D2.4 Каталог на услугите с граф на зависимостите и отговорните екипи (месец 4)\\*
\textit{Връзки} & Използва се от агента за зависимости (3.2), за оценката на каскадния риск (3.3) и за насочването на алертите (5.1).\\*
\textit{Екип (ч.м.)} & платформен инженер 0,5; SRE инженер 0,25; бизнес анализатор 0,25 (общо 1 ч.м.)\\*
\textit{Материални ресурси} & Няма нови разходи.\\*
\textit{Нематериални ресурси} & Достъп до трасировките и до системата за дежурства (PagerDuty/Opsgenie); време на собствениците на услугите (около 30 екипа по 1 час).\\*
\textit{Разходи} & персонал \textbf{5~600~\euro{}} (няма други преки разходи)\\
\midrule[0.3pt]
\multicolumn{2}{@{}p{\textwidth}@{}}{\textbf{Задача 2.5: Маскиране на лични и чувствителни данни}\hfill месец 3–4}\\*
\textit{Цел и същност} & Преди данните да стигнат до LLM или до хранилището, компонентът открива и заменя лични данни и тайни (имейли, IP адреси, токени, клиентски номера). Комбинираме регулярни изрази с библиотека с отворен код (напр. Microsoft Presidio) и проверяваме покритието върху извадка от реални логове.\\*
\textit{Резултат} & D2.5 Компонент за маскиране на чувствителни данни (месец 4)\\*
\textit{Връзки} & Прилага правилата от D1.4. Без него РП3 не може да работи с данни от prod; проверява се отново в 6.2.\\*
\textit{Екип (ч.м.)} & платформен инженер 0,25; специалист по сигурност 0,25; AI/LLM инженер 0,25 (общо 0,75 ч.м.)\\*
\textit{Материални ресурси} & Библиотека с отворен код, без лицензна такса.\\*
\textit{Нематериални ресурси} & Анонимизирана извадка от логове за тестване; консултации с екипа по сигурност.\\*
\textit{Разходи} & персонал \textbf{4~700~\euro{}} (няма други преки разходи)\\
\end{xltabular}}
{\small
\begin{xltabular}{\textwidth}{@{}L{2.55cm}X@{}}
\caption{Описание на работен пакет РП3}\\
\toprule
\multicolumn{2}{@{}p{\textwidth}@{}}{\textbf{Номер и име на работния пакет:} РП3: AI агенти за разследване}\\*
\multicolumn{2}{@{}l@{}}{\textbf{Начало (месец):} 3\qquad \textbf{Край (месец):} 6}\\*
\multicolumn{2}{@{}p{\textwidth}@{}}{\textbf{Цел:} (подцел 3) Автоматизиране на разследването на всеки алерт от специализирани AI агенти до хипотеза за първопричината}\\*
\multicolumn{2}{@{}p{\textwidth}@{}}{\textbf{Общо ресурси:} 6,5 ч.м.; 43~775~\euro{} (персонал 40~275~\euro{}, други преки разходи 3~500~\euro{})}\\
\midrule
\endfirsthead
\multicolumn{2}{@{}l@{}}{\textit{Таблица 2.\arabic{table}: РП3 (продължение)}}\\ \midrule
\endhead
\bottomrule
\endlastfoot
\multicolumn{2}{@{}p{\textwidth}@{}}{\textbf{Задача 3.1: Planning слой: намерение и декомпозиция на разследването}\hfill месец 3–4}\\*
\textit{Цел и същност} & Planning агентът чете нормализирания алерт и контекста му, определя какво трябва да се провери и разбива разследването на подзадачи: кои логове, кои health endpoint-и, кои зависими услуги. Използва LLM със структуриран изход (JSON), а промптовете се версионират като код.\\*
\textit{Резултат} & D3.1 Planning агент (месец 4)\\*
\textit{Връзки} & Използва D2.2 и интерфейсите от D1.5. Подзадачите му се изпълняват от агентите в 3.2.\\*
\textit{Екип (ч.м.)} & AI/LLM инженер 1; софтуерен архитект 0,25 (общо 1,25 ч.м.)\\*
\textit{Материални ресурси} & LLM API за разработка и експерименти (500~\euro{}).\\*
\textit{Нематериални ресурси} & Набор от реални алерти от dev и staging; знанието на SRE екипа как се разследва ръчно.\\*
\textit{Разходи} & персонал 8~250~\euro{} + други 500~\euro{} = \textbf{8~750~\euro{}}\\
\midrule[0.3pt]
\multicolumn{2}{@{}p{\textwidth}@{}}{\textbf{Задача 3.2: Execution слой: маршрутизиране и специализирани агенти}\hfill месец 4–5}\\*
\textit{Цел и същност} & Маршрутизаторът разпределя подзадачите към специализирани агенти за логове, за health endpoint-и и за зависимости. Всеки агент ползва конекторите от 2.3 като инструменти; подзадачите вървят паралелно с ограничение на времето и броя стъпки.\\*
\textit{Резултат} & D3.2 Специализирани агенти за логове, health endpoint-и и зависимости (месец 5)\\*
\textit{Връзки} & Изпълнява плана от 3.1 с конекторите (2.3) и графа (2.4). Резултатите отиват в 3.3; агентите могат да пускат проверките от 5.3.\\*
\textit{Екип (ч.м.)} & AI/LLM инженер 1,75; платформен инженер 0,5 (общо 2,25 ч.м.)\\*
\textit{Материални ресурси} & LLM API за разработка (1 500~\euro{}): агентите правят много последователни извиквания.\\*
\textit{Нематериални ресурси} & Достъпите от 2.3; примери за ръчни разследвания от минали инциденти.\\*
\textit{Разходи} & персонал 14~275~\euro{} + други 1~500~\euro{} = \textbf{15~775~\euro{}}\\
\midrule[0.3pt]
\multicolumn{2}{@{}p{\textwidth}@{}}{\textbf{Задача 3.3: Aggregation слой: обобщение и хипотеза за първопричина}\hfill месец 4–5}\\*
\textit{Цел и същност} & Обединяваме находките на агентите в доклад за инцидента: засегнати услуги, ключови редове от логовете, оценка на каскадния риск по графа и хипотеза за първопричината с ниво на увереност. Докладът се записва в споделения контекст.\\*
\textit{Резултат} & D3.3 Aggregation компонент с автоматичен доклад за инцидента (месец 5)\\*
\textit{Връзки} & Използва 3.2 и графа от 2.4. Докладът е вход за класификацията (4.2), тикетите (5.2) и докладите по NIS2 (5.5).\\*
\textit{Екип (ч.м.)} & AI/LLM инженер 0,75 (общо 0,75 ч.м.)\\*
\textit{Материални ресурси} & LLM API (500~\euro{}).\\*
\textit{Нематериални ресурси} & Шаблонът за postmortem на организацията като образец за структурата на доклада.\\*
\textit{Разходи} & персонал 4~875~\euro{} + други 500~\euro{} = \textbf{5~375~\euro{}}\\
\midrule[0.3pt]
\multicolumn{2}{@{}p{\textwidth}@{}}{\textbf{Задача 3.4: Разширяване на чата „Ask Alerts“}\hfill месец 6}\\*
\textit{Цел и същност} & Чатът от прототипа започва да работи с всички среди, вкл. prod, и да ползва агентите от 3.2 вместо собствена логика. Добавяме права по роли и история на разговорите, за да може инженерът да продължи разследването от там, където е спрял.\\*
\textit{Резултат} & D3.4 Нова версия на „Ask Alerts“ за всички среди (месец 6)\\*
\textit{Връзки} & Стъпва на 3.2 и 3.3. Показва се на демонстрацията в месец 6 (7.2); оценките на отговорите се ползват в 4.4.\\*
\textit{Екип (ч.м.)} & AI/LLM инженер 0,5; frontend разработчик 0,5 (общо 1 ч.м.)\\*
\textit{Материални ресурси} & LLM API (300~\euro{}).\\*
\textit{Нематериални ресурси} & Системата за единен вход (SSO) на организацията за правата по роли.\\*
\textit{Разходи} & персонал 5~750~\euro{} + други 300~\euro{} = \textbf{6~050~\euro{}}\\
\midrule[0.3pt]
\multicolumn{2}{@{}p{\textwidth}@{}}{\textbf{Задача 3.5: Оценка на агентите и защитни механизми (guardrails)}\hfill месец 5–6}\\*
\textit{Цел и същност} & От анотираните инциденти (4.1) правим тестов набор с известна първопричина и автоматично мерим дали хипотезата на агентите е вярна и за колко време идва. Добавяме защити: агентите имат само права за четене, текстът от логовете не може да променя инструкциите им (prompt injection), има лимит за време и цена на едно разследване.\\*
\textit{Резултат} & D3.5 Тестов набор и доклад за качеството на агентите (месец 6)\\*
\textit{Връзки} & Използва 4.1 и правилата от D1.4. Показва дали е постигнат MS3 (вярна хипотеза в поне 70\% от случаите, анализ до 5 минути); резултатите влизат в 6.2.\\*
\textit{Екип (ч.м.)} & AI/LLM инженер 0,5; QA инженер 0,5; специалист по сигурност 0,25 (общо 1,25 ч.м.)\\*
\textit{Материални ресурси} & LLM API за многократните тестови прогони (700~\euro{}).\\*
\textit{Нематериални ресурси} & Анотираните данни от 4.1; насоки от екипа по сигурност за атаки срещу LLM.\\*
\textit{Разходи} & персонал 7~125~\euro{} + други 700~\euro{} = \textbf{7~825~\euro{}}\\
\end{xltabular}}
{\small
\begin{xltabular}{\textwidth}{@{}L{2.55cm}X@{}}
\caption{Описание на работен пакет РП4}\\
\toprule
\multicolumn{2}{@{}p{\textwidth}@{}}{\textbf{Номер и име на работния пакет:} РП4: Класификация и самообучение}\\*
\multicolumn{2}{@{}l@{}}{\textbf{Начало (месец):} 3\qquad \textbf{Край (месец):} 7}\\*
\multicolumn{2}{@{}p{\textwidth}@{}}{\textbf{Цел:} (подцел 4) Надеждна автоматична класификация на инцидентите (фалшив алерт, деградация, прекъсване), която се подобрява с натрупването на история}\\*
\multicolumn{2}{@{}p{\textwidth}@{}}{\textbf{Общо ресурси:} 5,5 ч.м.; 34~300~\euro{} (персонал 33~100~\euro{}, други преки разходи 1~200~\euro{})}\\
\midrule
\endfirsthead
\multicolumn{2}{@{}l@{}}{\textit{Таблица 2.\arabic{table}: РП4 (продължение)}}\\ \midrule
\endhead
\bottomrule
\endlastfoot
\multicolumn{2}{@{}p{\textwidth}@{}}{\textbf{Задача 4.1: Събиране и анотиране на исторически инциденти}\hfill месец 3–4}\\*
\textit{Цел и същност} & Избираме около 1 500 алерта от последните 6–12 месеца, свързваме ги с тикетите и postmortem документите и ги привеждаме към формата от 2.1. Двама SRE инженери ги анотират независимо по таксономията от 1.3; съгласието им мерим с коефициента на Cohen (kappa), а спорните случаи решават заедно. Данните делим на обучаващи, валидационни и тестови.\\*
\textit{Резултат} & D4.1 Анотиран набор от данни за обучение и тестване (месец 4)\\*
\textit{Връзки} & Използва 1.3 и формата от 2.1. На него стъпват 3.5, 4.2 и 4.3.\\*
\textit{Екип (ч.м.)} & data scientist 0,75; SRE инженер 0,75 (общо 1,5 ч.м.)\\*
\textit{Материални ресурси} & Инструмент за анотиране с отворен код (Label Studio) в общата среда. Няма нови разходи.\\*
\textit{Нематериални ресурси} & Историческите данни в Splunk и тикет системата; експертизата на двамата анотатори.\\*
\textit{Разходи} & персонал \textbf{9~000~\euro{}} (няма други преки разходи)\\
\midrule[0.3pt]
\multicolumn{2}{@{}p{\textwidth}@{}}{\textbf{Задача 4.2: LLM класификация с вероятност за всяка категория}\hfill месец 4–5}\\*
\textit{Цел и същност} & Classification слоят подава на LLM доклада от 3.3 и критериите от таксономията и получава вероятност за всяка от трите категории. Върху валидационните данни проверяваме доколко вероятностите отговарят на реалната честота (калибровка).\\*
\textit{Резултат} & D4.2 Класификационен модул (Classification слой) (месец 5)\\*
\textit{Връзки} & Използва 3.3, 1.3 и 4.1. Комбинира се с класическите модели в 4.5.\\*
\textit{Екип (ч.м.)} & data scientist 0,75; AI/LLM инженер 0,25 (общо 1 ч.м.)\\*
\textit{Материални ресурси} & LLM API за прогоните върху целия набор (1 200~\euro{}).\\*
\textit{Нематериални ресурси} & Таксономията от D1.3.\\*
\textit{Разходи} & персонал 6~125~\euro{} + други 1~200~\euro{} = \textbf{7~325~\euro{}}\\
\midrule[0.3pt]
\multicolumn{2}{@{}p{\textwidth}@{}}{\textbf{Задача 4.3: Класически модели (логистична регресия, random forest)}\hfill месец 5–6}\\*
\textit{Цел и същност} & От алертите извличаме признаци (източник, услуга, час, честота на повторение, скорошни внедрявания, стойности на златните сигнали) и обучаваме логистична регресия и random forest с кръстосана проверка. Сравняваме ги с LLM класификацията.\\*
\textit{Резултат} & D4.3 Обучени модели и доклад от сравнението им (месец 6)\\*
\textit{Връзки} & Използва 4.1. Моделите се комбинират в 4.5 и се преобучават от механизма в 4.4.\\*
\textit{Екип (ч.м.)} & data scientist 1 (общо 1 ч.м.)\\*
\textit{Материални ресурси} & Процесорно време в общата среда (не е нужен GPU); библиотеки с отворен код (scikit-learn, MLflow).\\*
\textit{Нематериални ресурси} & Анотираните данни от 4.1.\\*
\textit{Разходи} & персонал \textbf{6~000~\euro{}} (няма други преки разходи)\\
\midrule[0.3pt]
\multicolumn{2}{@{}p{\textwidth}@{}}{\textbf{Задача 4.4: Обратна връзка от дежурните и периодично преобучение}\hfill месец 6–7}\\*
\textit{Цел и същност} & Дежурният потвърждава или поправя категорията от таблото или от тикета. Корекциите стават нови етикети, моделите се преобучават автоматично всяка седмица, а при промяна в разпределението на данните (data drift) се вдига сигнал.\\*
\textit{Резултат} & D4.4 Механизъм за обратна връзка и автоматично преобучение (месец 7)\\*
\textit{Връзки} & Преобучава моделите от 4.3; бутоните за обратна връзка са в таблото (5.4) и в тикетите (5.2). Това е Learning слоят от архитектурата.\\*
\textit{Екип (ч.м.)} & платформен инженер 0,75; data scientist 0,5 (общо 1,25 ч.м.)\\*
\textit{Материални ресурси} & Изчислителни ресурси в общата среда. Няма нови разходи.\\*
\textit{Нематериални ресурси} & Готовността на дежурните да оценяват класификацията (уговаря се в обучението, 6.4).\\*
\textit{Разходи} & персонал \textbf{7~350~\euro{}} (няма други преки разходи)\\
\midrule[0.3pt]
\multicolumn{2}{@{}p{\textwidth}@{}}{\textbf{Задача 4.5: Комбиниране на моделите и финална оценка}\hfill месец 7}\\*
\textit{Цел и същност} & Комбинираме LLM и класическите модели (претеглено или чрез stacking) и избираме праговете така, че да се разпознават поне 95\% от прекъсванията, дори с цената на повече ескалирани деградации. Оценката е върху тестови данни, които не са ползвани досега.\\*
\textit{Резултат} & D4.5 Финален класификатор и доклад от оценката (месец 7)\\*
\textit{Връзки} & Използва 4.2–4.4. Показва дали е постигнат MS4; класификаторът се ползва от 5.1 и се тества в 6.1.\\*
\textit{Екип (ч.м.)} & data scientist 0,5; AI/LLM инженер 0,25 (общо 0,75 ч.м.)\\*
\textit{Материални ресурси} & Няма нови разходи.\\*
\textit{Нематериални ресурси} & Тестовите данни от 4.1.\\*
\textit{Разходи} & персонал \textbf{4~625~\euro{}} (няма други преки разходи)\\
\end{xltabular}}
{\small
\begin{xltabular}{\textwidth}{@{}L{2.55cm}X@{}}
\caption{Описание на работен пакет РП5}\\
\toprule
\multicolumn{2}{@{}p{\textwidth}@{}}{\textbf{Номер и име на работния пакет:} РП5: Оперативна реакция и автоматизация}\\*
\multicolumn{2}{@{}l@{}}{\textbf{Начало (месец):} 5\qquad \textbf{Край (месец):} 7}\\*
\multicolumn{2}{@{}p{\textwidth}@{}}{\textbf{Цел:} (подцел 5) Автоматизиране на действията след класификацията: насочване към отговорния екип, регистриране на инцидента и автоматични проверки}\\*
\multicolumn{2}{@{}p{\textwidth}@{}}{\textbf{Общо ресурси:} 5,75 ч.м.; 33~375~\euro{} (персонал 32~375~\euro{}, други преки разходи 1~000~\euro{})}\\
\midrule
\endfirsthead
\multicolumn{2}{@{}l@{}}{\textit{Таблица 2.\arabic{table}: РП5 (продължение)}}\\ \midrule
\endhead
\bottomrule
\endlastfoot
\multicolumn{2}{@{}p{\textwidth}@{}}{\textbf{Задача 5.1: Насочване, ескалация и затваряне на алертите}\hfill месец 5–6}\\*
\textit{Цел и същност} & Правилата използват категорията и мапинга на компонентите: фалшивите алерти се затварят автоматично с обосновка и запис в одитния журнал, деградациите отиват в опашката на отговорния екип, а прекъсванията се ескалират към дежурния през PagerDuty/Opsgenie.\\*
\textit{Резултат} & D5.1 Модул за насочване, ескалация и затваряне на алертите (месец 6)\\*
\textit{Връзки} & Използва мапинга (2.4), класификацията (4.2, а по-късно 4.5) и правилата от 1.3. Действията му се регистрират чрез 5.2.\\*
\textit{Екип (ч.м.)} & DevOps/интеграционен инженер 1; SRE инженер 0,25 (общо 1,25 ч.м.)\\*
\textit{Материални ресурси} & Наличният лиценз за системата за дежурства. Няма нови разходи.\\*
\textit{Нематериални ресурси} & Политиката за ескалации на организацията; одобрение на правилата от SRE екипа.\\*
\textit{Разходи} & персонал \textbf{7~300~\euro{}} (няма други преки разходи)\\
\midrule[0.3pt]
\multicolumn{2}{@{}p{\textwidth}@{}}{\textbf{Задача 5.2: Интеграция със Splunk и тикет системата}\hfill месец 5–6}\\*
\textit{Цел и същност} & За всеки потвърден инцидент автоматично се създават събитие в Splunk и тикет с доклада от 3.3, линкове към логовете и отговорния екип. Ако за същия проблем вече има отворен тикет, новият алерт се добавя към него.\\*
\textit{Резултат} & D5.2 Интеграция със Splunk и тикет системата (месец 6)\\*
\textit{Връзки} & Използва 3.3 и 5.1. В тикета има бутони за обратна връзка (4.4); данните се ползват за докладите в 5.5.\\*
\textit{Екип (ч.м.)} & DevOps/интеграционен инженер 1 (общо 1 ч.м.)\\*
\textit{Материални ресурси} & Наличните лицензи за Splunk и тикет системата (обемът на новите събития е малък и не изисква нов лиценз).\\*
\textit{Нематериални ресурси} & Сервизни акаунти с права за запис; съгласуване на полетата на тикета с екипите.\\*
\textit{Разходи} & персонал \textbf{5~800~\euro{}} (няма други преки разходи)\\
\midrule[0.3pt]
\multicolumn{2}{@{}p{\textwidth}@{}}{\textbf{Задача 5.3: Автоматични проверки и действия (Playwright, n8n)}\hfill месец 5–7}\\*
\textit{Цел и същност} & Пишем Playwright сценарии, които минават през ключовите потребителски пътища и потвърждават дали услугата наистина е недостъпна, и n8n потоци за безопасни действия: събиране на диагностика, уведомяване и рестарт от одобрен списък след потвърждение от човек.\\*
\textit{Резултат} & D5.3 Библиотека от Playwright сценарии и n8n потоци (месец 7)\\*
\textit{Връзки} & Пускат се от агентите (3.2) и от правилата (5.1); тестват се в 6.1.\\*
\textit{Екип (ч.м.)} & DevOps/интеграционен инженер 1; QA инженер 0,25 (общо 1,25 ч.м.)\\*
\textit{Материални ресурси} & Инструменти с отворен код (Playwright, n8n Community Edition, хостван в общата среда). Няма нови разходи.\\*
\textit{Нематериални ресурси} & Тестови акаунти в услугите за синтетичните проверки; списък с одобрени автоматични действия от собствениците на услугите.\\*
\textit{Разходи} & персонал \textbf{6~925~\euro{}} (няма други преки разходи)\\
\midrule[0.3pt]
\multicolumn{2}{@{}p{\textwidth}@{}}{\textbf{Задача 5.4: Разширяване на таблото за наблюдение}\hfill месец 6–7}\\*
\textit{Цел и същност} & Към таблото от прототипа добавяме prod средата, резултата от класификацията с вероятностите, статуса на тикетите, бутоните за обратна връзка и показателите MTTD и MTTR в реално време.\\*
\textit{Резултат} & D5.4 Нова версия на таблото (месец 7)\\*
\textit{Връзки} & Чете данни от 2.2; съдържа обратната връзка за 4.4; показателите от него се ползват в 6.5.\\*
\textit{Екип (ч.м.)} & frontend разработчик 1; бизнес анализатор 0,25 (общо 1,25 ч.м.)\\*
\textit{Материални ресурси} & Няма нови разходи.\\*
\textit{Нематериални ресурси} & Кратки тестове за удобство с 3–5 дежурни инженери.\\*
\textit{Разходи} & персонал \textbf{6~200~\euro{}} (няма други преки разходи)\\
\midrule[0.3pt]
\multicolumn{2}{@{}p{\textwidth}@{}}{\textbf{Задача 5.5: Автоматични доклади за значими инциденти (NIS2)}\hfill месец 6–7}\\*
\textit{Цел и същност} & При инцидент, който може да е значим по NIS2 [4], системата подготвя чернова на ранното предупреждение (до 24 часа) и на уведомлението (до 72 часа) във формата на националния екип за реагиране (CSIRT). Черновата се изпраща само след одобрение от човек.\\*
\textit{Резултат} & D5.5 Генератор на доклади за ранно предупреждение по NIS2 (месец 7)\\*
\textit{Връзки} & Използва изискванията от 1.4, доклада от 3.3 и данните от 5.2.\\*
\textit{Екип (ч.м.)} & DevOps/интеграционен инженер 0,5; специалист по сигурност 0,25; AI/LLM инженер 0,25 (общо 1 ч.м.)\\*
\textit{Материални ресурси} & Правен преглед на шаблоните на докладите (1 000~\euro{}).\\*
\textit{Нематериални ресурси} & Указанията на националния CSIRT; процедурата на организацията за докладване.\\*
\textit{Разходи} & персонал 6~150~\euro{} + други 1~000~\euro{} = \textbf{7~150~\euro{}}\\
\end{xltabular}}
{\small
\begin{xltabular}{\textwidth}{@{}L{2.55cm}X@{}}
\caption{Описание на работен пакет РП6}\\
\toprule
\multicolumn{2}{@{}p{\textwidth}@{}}{\textbf{Номер и име на работния пакет:} РП6: Тестване, внедряване и оценка}\\*
\multicolumn{2}{@{}l@{}}{\textbf{Начало (месец):} 6\qquad \textbf{Край (месец):} 9}\\*
\multicolumn{2}{@{}p{\textwidth}@{}}{\textbf{Цел:} (подцел 6) Внедряване на системата в 24/7 експлоатация във всички среди и оценка на ефекта ѝ спрямо базовите показатели}\\*
\multicolumn{2}{@{}p{\textwidth}@{}}{\textbf{Общо ресурси:} 7,25 ч.м.; 60~950~\euro{} (персонал 41~550~\euro{}, други преки разходи 19~400~\euro{})}\\
\midrule
\endfirsthead
\multicolumn{2}{@{}l@{}}{\textit{Таблица 2.\arabic{table}: РП6 (продължение)}}\\ \midrule
\endhead
\bottomrule
\endlastfoot
\multicolumn{2}{@{}p{\textwidth}@{}}{\textbf{Задача 6.1: Тестове от край до край и тестове за натоварване}\hfill месец 6–7}\\*
\textit{Цел и същност} & В staging пускаме отново исторически инциденти и изкуствено предизвикани повреди (fault injection) и проверяваме целия път от алерта до тикета. С k6 натоварваме системата до 10 пъти над нормалния поток и проверяваме закъснение под 10 s при приемането и анализ до 5 минути.\\*
\textit{Резултат} & D6.1 Доклад от тестването (месец 7)\\*
\textit{Връзки} & Тества РП2–РП5 срещу критериите за приемане от 1.2. Успешният доклад е условие за внедряването в prod (6.3).\\*
\textit{Екип (ч.м.)} & QA инженер 1; платформен инженер 0,5 (общо 1,5 ч.м.)\\*
\textit{Материални ресурси} & Временна среда за тестовете за натоварване (1 000~\euro{}) и LLM API за тях (800~\euro{}); k6 (отворен код).\\*
\textit{Нематериални ресурси} & Сценариите от историческите инциденти; критериите за приемане от D1.2.\\*
\textit{Разходи} & персонал 7~400~\euro{} + други 1~800~\euro{} = \textbf{9~200~\euro{}}\\
\midrule[0.3pt]
\multicolumn{2}{@{}p{\textwidth}@{}}{\textbf{Задача 6.2: Проверка на сигурността и защитата на данните}\hfill месец 6–7}\\*
\textit{Цел и същност} & Външна фирма прави тест за проникване (penetration test). Вътрешно проверяваме правата за достъп, управлението на тайните, маскирането на данни и устойчивостта на агентите на prompt injection. Откритите проблеми се отстраняват преди prod.\\*
\textit{Резултат} & D6.2 Доклад от проверката за сигурност (месец 7)\\*
\textit{Връзки} & Проверява изпълнението на 1.4, 2.5 и 3.5. Условие за 6.3.\\*
\textit{Екип (ч.м.)} & специалист по сигурност 0,75 (общо 0,75 ч.м.)\\*
\textit{Материални ресурси} & Външен тест за проникване (8 000~\euro{}).\\*
\textit{Нематериални ресурси} & Политиката от D1.4; изолирана тестова среда и достъпи за външната фирма.\\*
\textit{Разходи} & персонал 4~875~\euro{} + други 8~000~\euro{} = \textbf{12~875~\euro{}}\\
\midrule[0.3pt]
\multicolumn{2}{@{}p{\textwidth}@{}}{\textbf{Задача 6.3: Внедряване във всички среди и CI/CD}\hfill месец 7–8}\\*
\textit{Цел и същност} & Описваме инфраструктурата като код (Terraform, Helm), настройваме CI/CD и внедряваме поетапно: dev, staging, int-prod и накрая prod. Самата YTKI също се наблюдава (метрики и алерти за собствените ѝ откази). Пишем ръководство за експлоатация (runbook).\\*
\textit{Резултат} & D6.3 Внедрена система и ръководство за експлоатация (runbook) (месец 8)\\*
\textit{Връзки} & Изисква успешни 6.1 и 6.2. С нея започва пилотът (6.5) и се постига MS6.\\*
\textit{Екип (ч.м.)} & DevOps/интеграционен инженер 1,5; платформен инженер 0,5 (общо 2 ч.м.)\\*
\textit{Материални ресурси} & Облачна инфраструктура за prod за месеци 8–9 (2~$\times$~1 500~\euro{}).\\*
\textit{Нематериални ресурси} & Одобрение по процеса за управление на промените (change management) за prod.\\*
\textit{Разходи} & персонал 11~600~\euro{} + други 3~000~\euro{} = \textbf{14~600~\euro{}}\\
\midrule[0.3pt]
\multicolumn{2}{@{}p{\textwidth}@{}}{\textbf{Задача 6.4: Обучение на дежурните екипи}\hfill месец 8}\\*
\textit{Цел и същност} & Провеждаме три обучения по около два часа за около 40 дежурни инженери: как да четат доклада, кога да се доверят на автоматичното затваряне и как да дават обратна връзка. Подготвяме кратки видеа и често задавани въпроси.\\*
\textit{Резултат} & D6.4 Обучителни материали за дежурните екипи (месец 8)\\*
\textit{Връзки} & Използва runbook-а от 6.3. Трябва да приключи преди пилота (6.5).\\*
\textit{Екип (ч.м.)} & бизнес анализатор 0,5; SRE инженер 0,25; AI/LLM инженер 0,25 (общо 1 ч.м.)\\*
\textit{Материални ресурси} & Зала и кетъринг (600~\euro{}).\\*
\textit{Нематериални ресурси} & Време на обучаваните (около 80 човекочаса).\\*
\textit{Разходи} & персонал 5~525~\euro{} + други 600~\euro{} = \textbf{6~125~\euro{}}\\
\midrule[0.3pt]
\multicolumn{2}{@{}p{\textwidth}@{}}{\textbf{Задача 6.5: Пилотна експлоатация и оценка на ефекта}\hfill месец 8–9}\\*
\textit{Цел и същност} & YTKI работи 24/7 поне шест седмици и обработва всички алерти. Измерваме MTTD, MTTR, дела на алертите за ръчна обработка и наличността на YTKI и ги сравняваме с базовите стойности от 1.1. Екипът осигурява засилена поддръжка (hypercare) и събира обратна връзка.\\*
\textit{Резултат} & D6.5 Доклад от пилота и оценка на ефекта спрямо базовите показатели (месец 9)\\*
\textit{Връзки} & Сравнява с D1.1. Резултатите се представят в 7.4 и са основа за приемането в 8.4.\\*
\textit{Екип (ч.м.)} & SRE инженер 0,5; data scientist 0,5; DevOps/интеграционен инженер 0,5; AI/LLM инженер 0,5 (общо 2 ч.м.)\\*
\textit{Материални ресурси} & LLM API при пълния поток алерти (около 3 000~\euro{} на месец, общо 6 000~\euro{}); prod инфраструктурата от 6.3.\\*
\textit{Нематериални ресурси} & Дежурните, които работят със системата и дават обратна връзка.\\*
\textit{Разходи} & персонал 12~150~\euro{} + други 6~000~\euro{} = \textbf{18~150~\euro{}}\\
\end{xltabular}}
{\small
\begin{xltabular}{\textwidth}{@{}L{2.55cm}X@{}}
\caption{Описание на работен пакет РП7}\\
\toprule
\multicolumn{2}{@{}p{\textwidth}@{}}{\textbf{Номер и име на работния пакет:} РП7: Комуникации}\\*
\multicolumn{2}{@{}l@{}}{\textbf{Начало (месец):} 1\qquad \textbf{Край (месец):} 9}\\*
\multicolumn{2}{@{}p{\textwidth}@{}}{\textbf{Цел:} Широко разпространение на информацията за проекта и резултатите му до заинтересованите лица}\\*
\multicolumn{2}{@{}p{\textwidth}@{}}{\textbf{Общо ресурси:} 2 ч.м.; 12~550~\euro{} (персонал 10~750~\euro{}, други преки разходи 1~800~\euro{})}\\
\midrule
\endfirsthead
\multicolumn{2}{@{}l@{}}{\textit{Таблица 2.\arabic{table}: РП7 (продължение)}}\\ \midrule
\endhead
\bottomrule
\endlastfoot
\multicolumn{2}{@{}p{\textwidth}@{}}{\textbf{Задача 7.1: Комуникационна стратегия и план}\hfill месец 1}\\*
\textit{Цел и същност} & Според анализа на заинтересованите лица от Домашна работа 1 определяме какво, на кого, по какъв канал и колко често съобщаваме: Slack канал на проекта, wiki, демонстрации, месечен бюлетин към ръководството.\\*
\textit{Резултат} & D7.1 Комуникационен план (месец 1)\\*
\textit{Връзки} & Използва анализа от Домашна работа 1 и плана от 8.1. Изпълнява се чрез 7.2–7.4.\\*
\textit{Екип (ч.м.)} & бизнес анализатор 0,5 (общо 0,5 ч.м.)\\*
\textit{Материални ресурси} & Наличните вътрешни канали (Slack, Confluence). Няма разходи.\\*
\textit{Нематериални ресурси} & Анализът на заинтересованите лица.\\*
\textit{Разходи} & персонал \textbf{2~400~\euro{}} (няма други преки разходи)\\
\midrule[0.3pt]
\multicolumn{2}{@{}p{\textwidth}@{}}{\textbf{Задача 7.2: Демонстрации и обратна връзка от потребителите}\hfill месец 3, 6, 9}\\*
\textit{Цел и същност} & В месеци 3, 6 и 9 показваме на потребителите какво работи: първо потока от данни, после агентите и класификацията, накрая пилота. Обратната връзка обобщаваме и я превръщаме в задачи в backlog-а.\\*
\textit{Резултат} & D7.2 Отчети от демонстрациите с обобщена обратна връзка (месец 3, 6, 9)\\*
\textit{Връзки} & Показва резултатите от РП2–РП6; обратната връзка се връща към изискванията (1.2) и към 4.4.\\*
\textit{Екип (ч.м.)} & ръководител на проекта 0,25; AI/LLM инженер 0,25; бизнес анализатор 0,25 (общо 0,75 ч.м.)\\*
\textit{Материални ресурси} & Кетъринг и материали за трите демонстрации (300~\euro{}).\\*
\textit{Нематериални ресурси} & Време на участниците (дежурни, собственици на услуги, ръководство).\\*
\textit{Разходи} & персонал 4~325~\euro{} + други 300~\euro{} = \textbf{4~625~\euro{}}\\
\midrule[0.3pt]
\multicolumn{2}{@{}p{\textwidth}@{}}{\textbf{Задача 7.3: Вътрешна страница на проекта (wiki)}\hfill месец 2–9}\\*
\textit{Цел и същност} & Създаваме страница в Confluence със статуса на проекта, архитектурата, документацията и често задаваните въпроси и я обновяваме след всеки спринт.\\*
\textit{Резултат} & D7.3 Страница на проекта във вътрешната wiki (месец 2)\\*
\textit{Връзки} & Тук се публикуват резултатите от всички пакети.\\*
\textit{Екип (ч.м.)} & бизнес анализатор 0,25 (общо 0,25 ч.м.)\\*
\textit{Материални ресурси} & Наличният лиценз за Confluence.\\*
\textit{Нематериални ресурси} & Няма.\\*
\textit{Разходи} & персонал \textbf{1~200~\euro{}} (няма други преки разходи)\\
\midrule[0.3pt]
\multicolumn{2}{@{}p{\textwidth}@{}}{\textbf{Задача 7.4: Представяне на резултатите}\hfill месец 9}\\*
\textit{Цел и същност} & Представяме резултатите от пилота пред ръководството и във вътрешна техническа лекция и подготвяме публикация (статия или доклад на конференция) за подхода, без поверителни данни.\\*
\textit{Резултат} & D7.4 Финална презентация и публикация (месец 9)\\*
\textit{Връзки} & Използва 6.5. Част от приключването на проекта (8.4).\\*
\textit{Екип (ч.м.)} & AI/LLM инженер 0,25; бизнес анализатор 0,25 (общо 0,5 ч.м.)\\*
\textit{Материални ресурси} & Такса за участие и пътни разходи за една конференция в България (1 500~\euro{}).\\*
\textit{Нематериални ресурси} & Одобрение от ръководството и екипа по сигурност за публикуваното съдържание.\\*
\textit{Разходи} & персонал 2~825~\euro{} + други 1~500~\euro{} = \textbf{4~325~\euro{}}\\
\end{xltabular}}
{\small
\begin{xltabular}{\textwidth}{@{}L{2.55cm}X@{}}
\caption{Описание на работен пакет РП8}\\
\toprule
\multicolumn{2}{@{}p{\textwidth}@{}}{\textbf{Номер и име на работния пакет:} РП8: Управление на проекта}\\*
\multicolumn{2}{@{}l@{}}{\textbf{Начало (месец):} 1\qquad \textbf{Край (месец):} 9}\\*
\multicolumn{2}{@{}p{\textwidth}@{}}{\textbf{Цел:} Качествено и в срок приключване на проекта}\\*
\multicolumn{2}{@{}p{\textwidth}@{}}{\textbf{Общо ресурси:} 4,25 ч.м.; 25~375~\euro{} (персонал 25~375~\euro{}, други преки разходи 0~\euro{})}\\
\midrule
\endfirsthead
\multicolumn{2}{@{}l@{}}{\textit{Таблица 2.\arabic{table}: РП8 (продължение)}}\\ \midrule
\endhead
\bottomrule
\endlastfoot
\multicolumn{2}{@{}p{\textwidth}@{}}{\textbf{Задача 8.1: Иницииране и планиране на проекта}\hfill месец 1}\\*
\textit{Цел и същност} & Провеждаме стартова среща със спонсорите, приемаме устава на проекта (project charter) и подготвяме плана: обхват, график, бюджет, роли и правила за работа на спринтове.\\*
\textit{Резултат} & D8.1 План за управление на проекта (месец 1)\\*
\textit{Връзки} & Рамка за всички останали задачи; заедно със 7.1 определя как се отчита напредъкът.\\*
\textit{Екип (ч.м.)} & ръководител на проекта 0,5 (общо 0,5 ч.м.)\\*
\textit{Материални ресурси} & Наличната Jira.\\*
\textit{Нематериални ресурси} & Решенията на спонсорите за обхвата и приоритетите.\\*
\textit{Разходи} & персонал \textbf{3~000~\euro{}} (няма други преки разходи)\\
\midrule[0.3pt]
\multicolumn{2}{@{}p{\textwidth}@{}}{\textbf{Задача 8.2: Управление на рисковете и качеството}\hfill месец 1–9}\\*
\textit{Цел и същност} & Водим регистър на рисковете (напр. забавен достъп до prod, грешни хипотези на LLM, по-високи от планираните разходи за API) с вероятност, влияние и мерки и го преглеждаме всеки месец. Планът за качеството определя Definition of Done, преглед на кода и минимално покритие с тестове.\\*
\textit{Резултат} & D8.2 Регистър на рисковете и план за качеството (месец 2)\\*
\textit{Връзки} & Покрива всички пакети; рискът за разходите е свързан с резерва в бюджета (раздел 3).\\*
\textit{Екип (ч.м.)} & ръководител на проекта 0,75; QA инженер 0,25 (общо 1 ч.м.)\\*
\textit{Материални ресурси} & Няма.\\*
\textit{Нематериални ресурси} & Опитът на организацията от предишни проекти.\\*
\textit{Разходи} & персонал \textbf{5~625~\euro{}} (няма други преки разходи)\\
\midrule[0.3pt]
\multicolumn{2}{@{}p{\textwidth}@{}}{\textbf{Задача 8.3: Наблюдение и контрол на изпълнението}\hfill месец 1–9}\\*
\textit{Цел и същност} & Планираме и преглеждаме всеки двуседмичен спринт, следим напредъка спрямо графика и разходите спрямо бюджета, веднъж месечно провеждаме управляващ комитет със спонсорите. В месеци 3 и 6 подготвяме междинни отчети.\\*
\textit{Резултат} & D8.3 Междинни отчети за напредъка (месец 3, 6)\\*
\textit{Връзки} & Следи всички задачи и междинните резултати MS1–MS6.\\*
\textit{Екип (ч.м.)} & ръководител на проекта 2 (общо 2 ч.м.)\\*
\textit{Материални ресурси} & Jira и табло за бюджета (налични).\\*
\textit{Нематериални ресурси} & Време на спонсорите за управляващия комитет (около 1 час месечно).\\*
\textit{Разходи} & персонал \textbf{12~000~\euro{}} (няма други преки разходи)\\
\midrule[0.3pt]
\multicolumn{2}{@{}p{\textwidth}@{}}{\textbf{Задача 8.4: Приключване на проекта}\hfill месец 9}\\*
\textit{Цел и същност} & Спонсорите приемат системата по индикаторите от Домашна работа 2; предаваме я на екипа по поддръжка, провеждаме среща за научените уроци (lessons learned) и пишем финалния отчет.\\*
\textit{Резултат} & D8.4 Финален отчет и протокол за приемане (месец 9)\\*
\textit{Връзки} & Използва D6.5 и D7.4.\\*
\textit{Екип (ч.м.)} & ръководител на проекта 0,5; софтуерен архитект 0,25 (общо 0,75 ч.м.)\\*
\textit{Материални ресурси} & Няма.\\*
\textit{Нематериални ресурси} & Екипът по поддръжка, който поема системата.\\*
\textit{Разходи} & персонал \textbf{4~750~\euro{}} (няма други преки разходи)\\
\end{xltabular}}
\renewcommand{\thetable}{\arabic{table}}\setcounter{table}{2}

\subsection*{2.3\quad Разпределение на задачите във времето}
Графикът (фиг.~3) показва задачите по двуседмични спринтове. Пакетите частично се застъпват, както е описано в методологията: всяка задача започва, когато задачите, от които зависи, са дали първия си използваем резултат. Ромбовете отбелязват междинните резултати MS1--MS6 от Домашна работа 2 в края на съответния месец.

\begin{figure}[p]
\centering
\begin{ganttchart}[
  x unit=0.56cm, y unit title=0.48cm, y unit chart=0.405cm,
  hgrid={draw=black!8}, vgrid={{draw=black!12,dotted},{draw=black!35}},
  title/.style={draw=black!40, fill=wpfill!35},
  title label font=\scriptsize\bfseries,
  bar/.style={draw=tkdraw, fill=tkfill}, bar height=0.55,
  bar label font=\scriptsize, bar label node/.append style={align=left},
  group/.style={draw=wpdraw, fill=wpfill}, group height=0.38, group peaks height=0.14, group peaks tip position=0,
  group label font=\scriptsize\bfseries,
  milestone/.append style={draw=msfill, fill=msfill}, milestone height=0.62,
  milestone label font=\scriptsize\itshape\color{msfill},
  canvas/.style={draw=black!40}
]{1}{18}
\gantttitle{Месец от началото на проекта}{18}\\
\gantttitlelist{1,...,9}{2}\\
\ganttgroup{РП1: Анализ и проектиране}{1}{4}\\
\ganttbar{1.1 Текущ процес и базови показатели}{1}{2}\\
\ganttbar{1.2 Изисквания}{3}{4}\\
\ganttbar{1.3 Таксономия на инцидентите}{2}{4}\\
\ganttbar{1.4 Сигурност и съответствие}{3}{4}\\
\ganttbar{1.5 Архитектура и интерфейси}{2}{4}\\
\ganttmilestone{MS1 Изисквания и архитектура}{4}\\
\ganttgroup{РП2: Приемане и разпространение на данните}{3}{8}\\
\ganttbar{2.1 Ingestion слой}{3}{6}\\
\ganttbar{2.2 Broadcasting слой}{4}{6}\\
\ganttbar{2.3 Конектори към източниците}{5}{8}\\
\ganttbar{2.4 Граф на зависимостите и мапинг}{4}{8}\\
\ganttbar{2.5 Маскиране на данни}{6}{8}\\
\ganttmilestone{MS2 Единен поток от данни}{8}\\
\ganttgroup{РП3: AI агенти за разследване}{5}{12}\\
\ganttbar{3.1 Planning слой}{5}{7}\\
\ganttbar{3.2 Execution слой и агенти}{7}{10}\\
\ganttbar{3.3 Aggregation слой}{8}{10}\\
\ganttbar{3.4 „Ask Alerts“}{11}{12}\\
\ganttbar{3.5 Оценка и guardrails}{9}{12}\\
\ganttmilestone{MS3 Автоматично разследване}{12}\\
\ganttgroup{РП4: Класификация и самообучение}{5}{14}\\
\ganttbar{4.1 Анотиране на исторически данни}{5}{8}\\
\ganttbar{4.2 LLM класификация}{8}{10}\\
\ganttbar{4.3 Класически модели}{9}{12}\\
\ganttbar{4.4 Обратна връзка и преобучение}{12}{14}\\
\ganttbar{4.5 Комбиниране и финална оценка}{13}{14}\\
\ganttmilestone{MS4 Надеждна класификация}{14}\\
\ganttgroup{РП5: Оперативна реакция и автоматизация}{9}{14}\\
\ganttbar{5.1 Насочване и ескалация}{9}{12}\\
\ganttbar{5.2 Splunk и тикет система}{9}{12}\\
\ganttbar{5.3 Playwright и n8n}{10}{14}\\
\ganttbar{5.4 Табло}{11}{14}\\
\ganttbar{5.5 Доклади по NIS2}{12}{14}\\
\ganttmilestone{MS5 Автоматизирана реакция}{14}\\
\ganttgroup{РП6: Тестване, внедряване и оценка}{12}{18}\\
\ganttbar{6.1 Тестове E2E и натоварване}{12}{14}\\
\ganttbar{6.2 Проверка на сигурността}{12}{14}\\
\ganttbar{6.3 Внедряване и CI/CD}{13}{16}\\
\ganttbar{6.4 Обучение на дежурните}{15}{16}\\
\ganttbar{6.5 Пилот и оценка на ефекта}{15}{18}\\
\ganttmilestone{MS6 Пилотна експлоатация}{16}\\
\ganttgroup{РП7: Комуникации}{1}{18}\\
\ganttbar{7.1 Комуникационен план}{1}{2}\\
\ganttbar{7.2 Демонстрации}{5}{6}\ganttbar{}{11}{12}\ganttbar{}{17}{18}\\
\ganttbar{7.3 Wiki на проекта}{3}{4}\ganttbar[bar/.append style={fill=mtfill, draw=tkdraw!60, dashed}]{}{5}{18}\\
\ganttbar{7.4 Представяне на резултатите}{17}{18}\\
\ganttgroup{РП8: Управление на проекта}{1}{18}\\
\ganttbar{8.1 Иницииране и планиране}{1}{2}\\
\ganttbar{8.2 Рискове и качество}{2}{18}\\
\ganttbar{8.3 Наблюдение и контрол}{1}{18}\\
\ganttbar{8.4 Приключване}{17}{18}\\
\end{ganttchart}

\vspace{4pt}
\begin{tikzpicture}[font=\scriptsize]
\filldraw[fill=wpfill, draw=wpdraw] (0,0) rectangle (0.6,0.18); \node[right] at (0.65,0.09) {работен пакет};
\filldraw[fill=tkfill, draw=tkdraw] (2.8,0) rectangle (3.4,0.18); \node[right] at (3.45,0.09) {задача};
\filldraw[fill=mtfill, draw=tkdraw!60, dashed] (4.6,0) rectangle (5.2,0.18); \node[right] at (5.25,0.09) {поддръжка (7.3)};
\fill[msfill] (8.0,0.09) -- (8.12,0.21) -- (8.24,0.09) -- (8.12,-0.03) -- cycle; \node[right] at (8.3,0.09) {междинен резултат (MS)};
\node[right] at (11.9,0.09) {колона = един спринт (2 седмици)};
\end{tikzpicture}
\caption{График на работата (Gantt chart) по спринтове и месеци.}
\label{fig:gantt}
\end{figure}
\clearpage


\subsection*{2.4\quad Списък на резултатите}
Таблица~3 изброява всички 38 резултата, подредени по месеца, в който се представят. Видът е: Д~--~документ (доклад, план, спецификация), С~--~софтуер (работещ компонент с документацията му), Дн~--~набор от данни, Др~--~друго (обучителни материали, wiki). Нивото на разпространение е: В~--~вътрешно (за организацията), О~--~ограничено (само за проектния екип и екипа по сигурност, защото съдържа лични данни или уязвимости), П~--~публично.
{\small
\begin{xltabular}{\textwidth}{@{}l X l c c c l@{}}
\caption{Списък на резултатите (deliverables)}\label{tab:deliv}\\
\toprule
\textbf{№} & \textbf{Име на резултата} & \textbf{РП} & \textbf{Вид} & \textbf{Ниво} & \textbf{Месец} & \textbf{Водеща роля}\\
\midrule
\endfirsthead
\multicolumn{7}{@{}l}{\textit{Таблица 3 (продължение)}}\\
\toprule
\textbf{№} & \textbf{Име на резултата} & \textbf{РП} & \textbf{Вид} & \textbf{Ниво} & \textbf{Месец} & \textbf{Водеща роля}\\
\midrule
\endhead
\bottomrule
\endlastfoot
D1.1 & Доклад за текущото състояние и базовите показатели & РП1 & Д & В & 1 & Бизнес анализатор\\
D7.1 & Комуникационен план & РП7 & Д & В & 1 & Бизнес анализатор\\
D8.1 & План за управление на проекта & РП8 & Д & В & 1 & Ръководител\\
D1.2 & Спецификация на изискванията & РП1 & Д & В & 2 & Бизнес анализатор\\
D1.3 & Таксономия на инцидентите с критерии и примери & РП1 & Д & В & 2 & SRE инж.\\
D1.4 & Политика за сигурност и защита на данните, вкл. правила за данните към LLM & РП1 & Д & О & 2 & Сигурност\\
D1.5 & Архитектурен документ на YTKI & РП1 & Д & В & 2 & Архитект\\
D7.3 & Страница на проекта във вътрешната wiki & РП7 & Др & В & 2 & Бизнес анализатор\\
D8.2 & Регистър на рисковете и план за качеството & РП8 & Д & В & 2 & Ръководител\\
D2.1 & Ingestion услуга с единен модел на алерта & РП2 & С & В & 3 & Платформен инж.\\
D2.2 & Broadcasting услуга със споделен контекст & РП2 & С & В & 3 & Платформен инж.\\
D7.2 & Отчети от демонстрациите с обобщена обратна връзка & РП7 & Д & В & 3, 6, 9 & Ръководител\\
D8.3 & Междинни отчети за напредъка & РП8 & Д & В & 3, 6 & Ръководител\\
D2.3 & Библиотека от конектори към източниците на данни & РП2 & С & В & 4 & Платформен инж.\\
D2.4 & Каталог на услугите с граф на зависимостите и отговорните екипи & РП2 & Дн & В & 4 & Платформен инж.\\
D2.5 & Компонент за маскиране на чувствителни данни & РП2 & С & В & 4 & Платформен инж.\\
D3.1 & Planning агент & РП3 & С & В & 4 & AI/LLM инж.\\
D4.1 & Анотиран набор от данни за обучение и тестване & РП4 & Дн & О & 4 & Data scientist\\
D3.2 & Специализирани агенти за логове, health endpoint-и и зависимости & РП3 & С & В & 5 & AI/LLM инж.\\
D3.3 & Aggregation компонент с автоматичен доклад за инцидента & РП3 & С & В & 5 & AI/LLM инж.\\
D4.2 & Класификационен модул (Classification слой) & РП4 & С & В & 5 & Data scientist\\
D3.4 & Нова версия на „Ask Alerts“ за всички среди & РП3 & С & В & 6 & AI/LLM инж.\\
D3.5 & Тестов набор и доклад за качеството на агентите & РП3 & Д & В & 6 & AI/LLM инж.\\
D4.3 & Обучени модели и доклад от сравнението им & РП4 & С & В & 6 & Data scientist\\
D5.1 & Модул за насочване, ескалация и затваряне на алертите & РП5 & С & В & 6 & DevOps инж.\\
D5.2 & Интеграция със Splunk и тикет системата & РП5 & С & В & 6 & DevOps инж.\\
D4.4 & Механизъм за обратна връзка и автоматично преобучение & РП4 & С & В & 7 & Платформен инж.\\
D4.5 & Финален класификатор и доклад от оценката & РП4 & С & В & 7 & Data scientist\\
D5.3 & Библиотека от Playwright сценарии и n8n потоци & РП5 & С & В & 7 & DevOps инж.\\
D5.4 & Нова версия на таблото & РП5 & С & В & 7 & Frontend разр.\\
D5.5 & Генератор на доклади за ранно предупреждение по NIS2 & РП5 & С & В & 7 & DevOps инж.\\
D6.1 & Доклад от тестването & РП6 & Д & В & 7 & QA инж.\\
D6.2 & Доклад от проверката за сигурност & РП6 & Д & О & 7 & Сигурност\\
D6.3 & Внедрена система и ръководство за експлоатация (runbook) & РП6 & С & В & 8 & DevOps инж.\\
D6.4 & Обучителни материали за дежурните екипи & РП6 & Др & В & 8 & Бизнес анализатор\\
D6.5 & Доклад от пилота и оценка на ефекта спрямо базовите показатели & РП6 & Д & В & 9 & SRE инж.\\
D7.4 & Финална презентация и публикация & РП7 & Д & П & 9 & AI/LLM инж.\\
D8.4 & Финален отчет и протокол за приемане & РП8 & Д & В & 9 & Ръководител\\
\end{xltabular}}

\section*{3\quad Необходими ресурси и бюджет}
\subsection*{3.1\quad Подход и допускания}
Бюджетът е изчислен отдолу нагоре (bottom-up estimation) [3]: за всяка задача в таблици 2.1--2.8 оценихме колко човекомесеца от всяка роля са нужни и какви други преки разходи има, а после сумирахме по пакети и категории. Сумите са в евро, защото от 1 януари 2026~г. еврото е националната валута на България [1].
\begin{itemize}
\item \textbf{Персонал.} Месечната ставка на роля е пълният разход на работодателя (брутна заплата и осигуровки за сметка на работодателя) за специалист със средно до високо ниво на опит. Ставките (4~500--7~000~\euro{}) съответстват на пазарните нива в България през 2026~г. [2]. Един човекомесец е около 21 работни дни.
\item \textbf{LLM.} Разходът е оценен по броя извиквания: при пилота очакваме около 100 алерта на ден след отсяването на дубликатите, около 10 извиквания на LLM за едно разследване и около 1~\euro{} на разследване, т.е. около 3~000~\euro{} на месец. В разработката разходът е по-малък и е разпределен по задачите, които го пораждат.
\item \textbf{Налични ресурси.} Организацията вече плаща за Splunk, Jira, Confluence, системата за дежурства, Slack и SSO, а екипът работи на собственото си оборудване и в офиса на организацията. Тези ресурси не влизат в бюджета на проекта, но са условие за него (раздел~3.3).
\item \textbf{Резерв.} Добавяме 10\% резерв за непредвидени разходи. Той покрива основните рискове от регистъра (8.2): по-високи разходи за LLM API, забавен достъп до prod и допълнителен спринт за корекции след пилота. Резервът се използва само с одобрение на управляващия комитет.
\end{itemize}


\subsection*{3.2\quad Човешки ресурси}
Таблица~4 описва ролите в екипа, а табл.~5 показва колко човекомесеца влага всяка роля във всеки пакет. Най-много работа има за AI/LLM инженерите, платформените и DevOps инженерите, защото те изграждат слоевете на системата. SRE инженерът е от дежурните екипи: той носи знанието за реалните инциденти и е главният представител на потребителите в екипа.
{\small
\begin{xltabular}{\textwidth}{@{}L{3.6cm} c X R{1.6cm}@{}}
\caption{Роли в екипа на проекта}\label{tab:team}\\
\toprule
\textbf{Роля} & \textbf{Хора} & \textbf{Основна отговорност} & \textbf{\euro{}/ч.м.}\\
\midrule
\endfirsthead
\toprule
\textbf{Роля} & \textbf{Хора} & \textbf{Основна отговорност} & \textbf{\euro{}/ч.м.}\\
\midrule
\endhead
\bottomrule
\endlastfoot
Ръководител на проекта & 1 & планиране, наблюдение и контрол, рискове, връзка със спонсорите (РП8) & 6~000\\
Бизнес анализатор & 1 & базови показатели, изисквания, комуникации, обучителни материали & 4~800\\
Софтуерен архитект & 1 & архитектура, интерфейси между агентите, технически решения & 7~000\\
Платформен инженер / инженер по данни & 2 & поток от данни, конектори, граф на зависимостите, маскиране, преобучение & 5~800\\
AI/LLM инженер & 2 & LLM агенти, промптове, „Ask Alerts“, защитни механизми & 6~500\\
Data scientist & 1 & анализ на данните, анотиране, класификационни модели и оценката им & 6~000\\
DevOps инженер / инженер по интеграции & 2 & насочване, интеграции със Splunk и тикет системата, n8n, CI/CD и внедряване & 5~800\\
Frontend разработчик & 1 & интерфейс на „Ask Alerts“ и таблото & 5~000\\
QA инженер & 1 & тестове на агентите, от край до край и за натоварване & 4~500\\
Специалист по сигурност & 1 & DPIA, политика за данните, NIS2, проверка за сигурност & 6~500\\
SRE инженер (от дежурните екипи) & 1 & таксономия, анотиране, правила за ескалация, пилот & 6~000\\
\end{xltabular}}
{\small\setlength{\tabcolsep}{4.2pt}
\begin{table}[H]\centering
\caption{Човекомесеци по роли и работни пакети и разход за персонал}\label{tab:pm}
\begin{tabular}{@{}lrrrrrrrrrr@{}}
\toprule
\textbf{Роля} & \textbf{РП1} & \textbf{РП2} & \textbf{РП3} & \textbf{РП4} & \textbf{РП5} & \textbf{РП6} & \textbf{РП7} & \textbf{РП8} & \textbf{Общо} & \textbf{\euro{}}\\
\midrule
Ръководител & 0,25 & -- & -- & -- & -- & -- & 0,25 & 3,75 & 4,25 & 25~500\\
Бизнес анализатор & 1,25 & 0,25 & -- & -- & 0,25 & 0,5 & 1,25 & -- & 3,5 & 16~800\\
Архитект & 1,25 & -- & 0,25 & -- & -- & -- & -- & 0,25 & 1,75 & 12~250\\
Платформен инж. & 0,25 & 3,5 & 0,5 & 0,75 & -- & 1 & -- & -- & 6 & 34~800\\
AI/LLM инж. & 0,25 & 0,25 & 4,5 & 0,5 & 0,25 & 0,75 & 0,5 & -- & 7 & 45~500\\
Data scientist & 0,75 & -- & -- & 3,5 & -- & 0,5 & -- & -- & 4,75 & 28~500\\
DevOps инж. & -- & 1,25 & -- & -- & 3,5 & 2 & -- & -- & 6,75 & 39~150\\
Frontend разр. & -- & -- & 0,5 & -- & 1 & -- & -- & -- & 1,5 & 7~500\\
QA инж. & -- & -- & 0,5 & -- & 0,25 & 1 & -- & 0,25 & 2 & 9~000\\
Сигурност & 0,75 & 0,25 & 0,25 & -- & 0,25 & 0,75 & -- & -- & 2,25 & 14~625\\
SRE инж. & 1 & 0,25 & -- & 0,75 & 0,25 & 0,75 & -- & -- & 3 & 18~000\\
\midrule
\textbf{Общо ч.м.} & \textbf{5,75} & \textbf{5,75} & \textbf{6,5} & \textbf{5,5} & \textbf{5,75} & \textbf{7,25} & \textbf{2} & \textbf{4,25} & \textbf{42,75} & \textbf{251~625}\\
\textbf{Персонал, \euro{}} & 34~700 & 33~500 & 40~275 & 33~100 & 32~375 & 41~550 & 10~750 & 25~375 & & \\
\bottomrule\end{tabular}\end{table}}

\subsection*{3.3\quad Материални и нематериални ресурси}
Освен хората проектът има нужда от облачна инфраструктура, от достъп до LLM и от няколко външни услуги. Таблица~6 изброява всички преки разходи, различни от персонала, и задачите, за които са нужни. Почти всички инструменти, които добавяме (n8n, Playwright, k6, Presidio, Label Studio, scikit-learn, MLflow), са с отворен код и не изискват лицензи.
{\small
\begin{table}[H]\centering
\caption{Други преки разходи (материални ресурси и външни услуги)}\label{tab:other}
\begin{tabularx}{\textwidth}{@{}l X l R{1.7cm}@{}}
\toprule
\textbf{Категория} & \textbf{Ресурс} & \textbf{Задача} & \textbf{\euro{}}\\
\midrule
Облачна инфраструктура & Изолирана среда за проби на архитектурата & 1.5 & 500\\
 & Обща dev/staging среда (Kubernetes, опашка, база), м. 2--9 & 2.1 & 7~200\\
 & Временна среда за тестове за натоварване & 6.1 & 1~000\\
 & Prod инфраструктура, м. 8--9 & 6.3 & 3~000\\
\addlinespace[2pt]
Използване на LLM (API) & Сравнение на LLM модели & 1.5 & 300\\
 & Разработка на агентите и чата & 3.1--3.4 & 2~800\\
 & Тестови прогони за оценка на агентите & 3.5 & 700\\
 & Класификация на анотирания набор & 4.2 & 1~200\\
 & Тестове за натоварване & 6.1 & 800\\
 & Пилотна експлоатация при пълен поток (2 месеца) & 6.5 & 6~000\\
\addlinespace[2pt]
Външни услуги & Правна консултация по GDPR и NIS2 & 1.4 & 2~000\\
 & Правен преглед на шаблоните за NIS2 & 5.5 & 1~000\\
 & Външен тест за проникване & 6.2 & 8~000\\
\addlinespace[2pt]
Обучения и събития & Работни срещи за изискванията & 1.2 & 300\\
 & Обучения на дежурните екипи & 6.4 & 600\\
 & Демонстрации за потребителите & 7.2 & 300\\
 & Участие в конференция & 7.4 & 1~500\\
\midrule
\textbf{Общо} & & & \textbf{37~200}\\
\bottomrule\end{tabularx}\end{table}}

Нематериалните ресурси не струват пари на проекта, но без тях задачите не могат да се изпълнят, затова ги планираме заедно с ръководителите на съответните екипи:
\begin{itemize}
\item \textbf{Достъпи и одобрения:} права само за четене до логовете, метриките и трасировките във всичките четири среди (2.3), права за запис в Splunk и тикет системата (5.2), одобрение на промените в prod (6.3).
\item \textbf{Данни:} историята на алертите, инцидентите и postmortem документите за последните 6--12 месеца (1.1, 4.1), дефинициите на SLO, каталогът на услугите и дежурните ротации (2.4).
\item \textbf{Знание:} съществуващият прототип (табло, граф, „Ask Alerts“), опитът на SRE екипа в разследването на инциденти и указанията на DPO и на екипа по сигурност.
\item \textbf{Време на хора извън екипа:} интервюта, работни срещи, потвърждаване на собствеността на услугите, обучения, демонстрации и управляващ комитет, общо около 250 часа (около 1,5 ч.м.), разпределени между много хора.
\item \textbf{Налични лицензи и оборудване:} Splunk, Jira, Confluence, PagerDuty/Opsgenie, Slack, SSO, лаптопи и офис на организацията.
\end{itemize}


\subsection*{3.4\quad Обобщен бюджет}
Таблица~7 събира разходите по категории и работни пакети. Най-скъп е РП6, защото включва теста за проникване, prod инфраструктурата и разходите за LLM при пилота. Персоналът е около 87\% от преките разходи, което е обичайно за софтуерен проект.
{\small\setlength{\tabcolsep}{3.4pt}
\begin{table}[H]\centering
\caption{Обобщен бюджет по категории и работни пакети (\euro{})}\label{tab:budget}
\begin{tabular}{@{}lrrrrrrrrr@{}}
\toprule
\textbf{Категория} & \textbf{РП1} & \textbf{РП2} & \textbf{РП3} & \textbf{РП4} & \textbf{РП5} & \textbf{РП6} & \textbf{РП7} & \textbf{РП8} & \textbf{Общо}\\
\midrule
Персонал & 34~700 & 33~500 & 40~275 & 33~100 & 32~375 & 41~550 & 10~750 & 25~375 & 251~625\\
Облачна инфраструктура & 500 & 7~200 & -- & -- & -- & 4~000 & -- & -- & 11~700\\
LLM API & 300 & -- & 3~500 & 1~200 & -- & 6~800 & -- & -- & 11~800\\
Външни услуги & 2~000 & -- & -- & -- & 1~000 & 8~000 & -- & -- & 11~000\\
Обучения и събития & 300 & -- & -- & -- & -- & 600 & 1~800 & -- & 2~700\\
\midrule
\textbf{Преки разходи} & \textbf{37~800} & \textbf{40~700} & \textbf{43~775} & \textbf{34~300} & \textbf{33~375} & \textbf{60~950} & \textbf{12~550} & \textbf{25~375} & \textbf{288~825}\\
Дял & 13\% & 14\% & 15\% & 12\% & 12\% & 21\% & 4\% & 9\% & 100\%\\
Резерв (10\%) & \multicolumn{8}{c}{\textit{управлява се на ниво проект (8.2, 8.3)}} & 28~882\\
\midrule
\textbf{Общ бюджет} & \multicolumn{8}{c}{} & \textbf{317~708}\\
\bottomrule\end{tabular}\end{table}}

Фигура~4 показва как се разпределят преките разходи по месеци. Разходите са сравнително равномерни в месеци 2--5, а пикът е в месеци 6--7, когато едновременно вървят РП3, РП4, РП5 и РП6. В месеци 8--9 екипът намалява, но остават разходите за prod инфраструктурата и за LLM при пилота. Кумулативната крива е базовата линия, спрямо която в 8.3 се следят реалните разходи.

\begin{figure}[H]
\centering
\begin{tikzpicture}
\begin{axis}[
  width=0.92\textwidth, height=5.6cm, ybar, bar width=0.55cm,
  xmin=0.4, xmax=9.6, xtick={1,...,9}, xlabel={Месец}, ylabel={хил. \euro{} на месец},
  ymin=0, ymax=55, axis y line*=left, axis x line*=bottom,
  tick label style={font=\scriptsize}, label style={font=\scriptsize},
  nodes near coords, nodes near coords style={font=\tiny, /pgf/number format/fixed, /pgf/number format/precision=1, /pgf/number format/use comma},
  /pgf/number format/use comma,
  ymajorgrids, grid style={black!10}]
\addplot[fill=wpfill, draw=wpdraw] coordinates {(1,20.7) (2,34.6) (3,31.2) (4,32.7) (5,34.0) (6,44.0) (7,43.4) (8,25.5) (9,22.7)};
\end{axis}
\begin{axis}[
  width=0.92\textwidth, height=5.6cm, xmin=0.4, xmax=9.6, ymin=0, ymax=300,
  axis y line*=right, axis x line=none, ylabel={кумулативно, хил. \euro{}},
  tick label style={font=\scriptsize}, label style={font=\scriptsize}, /pgf/number format/use comma]
\addplot[thick, color=msfill, mark=*, mark size=1.4pt] coordinates {(1,20.7) (2,55.3) (3,86.5) (4,119.2) (5,153.2) (6,197.1) (7,240.5) (8,266.1) (9,288.8)};
\end{axis}
\end{tikzpicture}
\caption{Преки разходи по месеци (стълбове) и кумулативно (линия), без резерва.}
\label{fig:cost}
\end{figure}

\textbf{Обобщение.} Проектът изисква 42,75 човекомесеца работа на екип от 14 души в продължение на 9 месеца. Преките разходи са 288~825~\euro{} (персонал 251~625~\euro{} и други 37~200~\euro{}), а с резерва от 10\% общият бюджет е \textbf{317~708~\euro{}}. За сравнение, при ориентировъчната цена на престоя от Домашна работа 1 (около 5~600 долара на минута) целият бюджет е приблизително колкото един час престой, т.е. проектът се изплаща, ако общо спести около час престой.


\section*{Литература}
\begin{enumerate}[label={[\arabic*]},leftmargin=2.2em,itemsep=2pt]
\item Европейска комисия, \textit{България приема еврото от 1 януари 2026 г.}, Access2Markets, 2026.
\item Lemon.io, \textit{Software developer rates in Bulgaria}, данни за Q3 2026.
\item Project Management Institute, \textit{A Guide to the Project Management Body of Knowledge (PMBOK Guide)}, 7th ed., PMI, 2021.
\item Директива (ЕС) 2022/2555 на Европейския парламент и на Съвета (NIS2), ОВ L 333, 27.12.2022.
\end{enumerate}
\end{document}
